\section*{الفصل \ref{ch:b2:enolates-aldol} --- الإينولات وأيونات الإينولات والتفاعل الألدولي}

\begin{solution}{exo:b2:enolates-aldol:1}
البيوتانون: البوت-1-ين-2-ول \ce{CH2=C(OH)CH2CH3} والبوت-2-ين-2-ول \ce{CH3C(OH)=CHCH3} ((\textit{E}) 
و(\textit{Z})). والهكسانون الحلقي: الهكس-1-ين حلقي-1-ول.
\end{solution}

\begin{solution}{exo:b2:enolates-aldol:2}
\babelsublr{\foreignlanguage{arabic}{الإيثان}~$<$~\foreignlanguage{arabic}{البروبانون}~$<$~\foreignlanguage{arabic}{3-أوكسوبيوتانوات الإيثيل (10.7)}~$<$~\foreignlanguage{arabic}{البنتان-2،4-ديون (9.0)}}.
\end{solution}

\begin{solution}{exo:b2:enolates-aldol:3}
3-هيدروكسي بيوتانال، ثم البوت-2-ينال بالتسخين.
\end{solution}

\begin{solution}{exo:b2:enolates-aldol:4}
بروبانديوات ثنائي الإيثيل، وإيثوكسيد الصوديوم، و1-بروموبيوتان؛ ثم الحلمهة إلى حمض البيوتيل بروبانديويك؛
ثم يزيل التسخين \ce{CO2}: \ce{CH3(CH2)3CH2COOH}، حمض الهكسانويك.
\end{solution}

\begin{solution}{exo:b2:enolates-aldol:5}
الحركي: نزع البروتون عند C6 (جهة \ce{CH2})، بواسطة LDA عند \qty{-78}{\degreeCelsius}. والترموديناميكي:
الرابطة المزدوجة الأكثر استبدالًا C1=C2، نحو الكربون الحامل للميثيل، بألكوكسيد في
كحوله في درجة حرارة الغرفة.
\end{solution}

\begin{solution}{exo:b2:enolates-aldol:6}
4-فينيل بوت-3-ين-2-ون (بنزال الأسيتون)، \ce{C6H5CH=CHCOCH3}. ليس للبنزالدهيد هيدروجين $\alpha$: فلا
يستطيع تكوين أيون إينولات ولا يكون إلا المحب للإلكترونات.
\end{solution}

\begin{solution}{exo:b2:enolates-aldol:7}
يهاجم أيون إينولات إحدى مجموعتي الميثيل (C1) الكربونيلَ الآخر (C5): حلقة خماسية،
هي 3-ميثيل بنت-2-ين حلقي-1-ون بعد نزع الماء. والبديل، أي \omterm{def:b2:enolates-aldol:enolate}{أيون الإينولات} عند C3 يهاجم C5،
كان سيكوّن حلقة ثلاثية متوترة.
\end{solution}

\begin{solution}{exo:b2:enolates-aldol:8}
\ce{2CH3COOC2H5 -> CH3COCH2COOC2H5 + C2H5OH} بوجود إيثوكسيد الصوديوم؛ والناتج، المنزوع البروتون،
يُستعاد على هيئة 3-أوكسوبيوتانوات الإيثيل عند التحميض.
\end{solution}

\begin{solution}{exo:b2:enolates-aldol:9}
البروبانون والإيثانال والأسيتوفينون (لها مجموعة \ce{CH3-CO}؛ فالإيثانال هو \ce{CH3CHO}).
أما البنتان-3-ون فلا.
\end{solution}

\begin{solution}{exo:b2:enolates-aldol:10}
يهاجم \omterm{def:b2:enolates-aldol:enolate}{أيون الإينولات} عند أحد الطرفين الإسترَ عند الطرف الآخر: 2-أوكسو بنتان حلقي-1-كربوكسيلات الإيثيل،
حلقة خماسية.
\end{solution}

\begin{solution}{exo:b2:enolates-aldol:11}
في 3-ميثيل بنتان-2-ون يكون المركز الفراغي (C3) هو الكربون $\alpha$: فالتحول إلى \omterm{def:b2:enolates-aldol:tautomerism}{إينول} يجعله مستويًا، 
وإعادة البرتنة من أي من الوجهين تحوّله إلى راسيمي. وفي 4-ميثيل هكسان-2-ون يقع المركز الفراغي عند C4،
في الموضع $\beta$ بالنسبة إلى الكربونيل، فلا يمسّه التحول إلى \omterm{def:b2:enolates-aldol:tautomerism}{إينول}.
\end{solution}

\begin{solution}{exo:b2:enolates-aldol:12}
البروبانون مع مكافئين من البنزالدهيد في وسط قاعدي. ومع ألدهيدين مختلفين، كثّف
البروبانون مع أحد الألدهيدين أولًا (مكافئ واحد، مع عزل بنزال الأسيتون)، ثم كثّف مجموعة
الميثيل المتبقية فيه مع الثاني. أما في وعاء واحد فيتنافس الألدهيدان ويعطيان خليطًا من
ثلاثة دينونات.
\end{solution}

\begin{solution}{pb:b2:enolates-aldol:1}
\textbf{1.} البروبانون، وله ستة هيدروجينات $\alpha$ (ثلاثة على كل ميثيل).
\textbf{2.} $\mathrm{CH_3{-}CO{-}CH_2^-} \leftrightarrow \mathrm{CH_3{-}C(O^-){=}CH_2}$.
\textbf{3.} ليس له هيدروجين على الكربون المجاور لكربونيله (فذلك الكربون ينتمي إلى الحلقة ولا
يحمل H).
\textbf{4.} البنزالدهيد، وهو ألدهيد بفائض، هو المحب للإلكترونات الأفضل؛ فيلقاه أيون إينولات البروبانون
أكثر مما يلقى كربونيل بروبانون.
\textbf{5.} للكيتون مجموعتان كربونيتان تدفعان الكثافة الإلكترونية نحو كربون كربونيله 
وتعيقان الاقتراب؛ وللألدهيد مجموعة واحدة.
\textbf{6.} لا يلزم إلا جزء صغير من \omterm{def:b2:enolates-aldol:enolate}{أيون الإينولات}: فهو يتجدد كلما تفاعل.
\textbf{7.} يرتبط كربون \omterm{def:b2:enolates-aldol:enolate}{أيون الإينولات} بكربون كربونيل البنزالدهيد: فيتكوّن ألكوكسيد
4-هيدروكسي-4-فينيل بيوتان-2-ون.
\textbf{8.} نزع هيدروجين $\alpha$ وفقد الهيدروكسيد من الكربون $\beta$: 4-فينيل بوت-3-ين-2-ون.
\textbf{9.} الرابطة المزدوجة المتكوّنة مترافقة مع الحلقة ومع الكربونيل كليهما.
\textbf{10.} لمجموعة الميثيل الأخرى ثلاثة هيدروجينات $\alpha$ بعد.
\textbf{11.} \ce{2C6H5CHO + CH3COCH3 -> C6H5CH=CHCOCH=CHC6H5 + 2H2O}.
\textbf{12.} نزعا الماء، وترسب الناتج، الذي يغادر المحلول.
\textbf{13.} اثنتان، كلتاهما (\textit{E}).
\textbf{14.} المتماكبات (\textit{Z}) تضع مجموعة الفينيل بجوار مجموعة الكربونيل: إجهاد فراغي.
\textbf{15.} رابطتا \ce{C=C}، ورابطة \ce{C=O} واحدة والحلقتان: نظام مترافق يمتد عبر الجزيء.
\textbf{16.} ترافقه الممتد يضيّق فجوة $\pi$ (\cref{prop:b2:huckel:gap}): فيمتص عند الطرف
الأزرق من المرئي ويبدو أصفر.
\textbf{17.} لا: فهو مستوٍ (أو يكاد) وله مستوي مرآة.
\textbf{18.} 106.12 و58.08 و\qty{234.30}{g/mol}.
\textbf{19.} البنزالدهيد $2.1 \times 1.045/106.12 = \qty{20.7}{mmol}$؛ والبروبانون $0.75 \times
0.7845/58.08 = \qty{10.1}{mmol}$.
\textbf{20.} يحتاج البروبانون إلى \qty{20.3}{mmol} من البنزالدهيد، والمتاح \qty{20.7}{mmol}:
فالبروبانون هو المحدد، بفارق ضئيل.
\textbf{21.} $0.01013 \times 234.30 = \qty{2.37}{g}$.
\textbf{22.} يبقى البنزالدهيد الفائض في الإيثانول ويُزال بالغسل.
\textbf{23.} ناتج التكاثف الأحادي، بنزال الأسيتون.
\textbf{24.} $0.75 \times 2.37 = \boldsymbol{\approx \qty{1.78}{g}}$ من ثنائي بنزال الأسيتون.
\end{solution}
